\section{Labelled rays and integration of PWS tensors}
\label{r15ray:section}

\paragraph{Status.}
The assertions in this section concern PWS tensors in finite ordinary
unit-moment presentations.  They do not assert a pushforward for arbitrary
actual forms or for arbitrary proper analytic maps.  The realization and
cutoff hypotheses in Proposition~\ref{r15ray:analytic} are part of its
statement.  Their verification for the Green operators is separate.

We use the additive valuation $v=-\log|\cdot|$.  All radii below are real;
no membership in the value group is imposed.  For a polyhedral set $P$,
$\mathrm{PS}^{p,q}(P)$ denotes compatible smooth Lagerberg tensors on a
finite polyhedral subdivision, including tangential agreement on faces.
Pullback under a piecewise affine map is used only in this PWS category.
We never use such a pullback for a delta-form of positive delta codimension.

\begin{definition}[A finite distance family]
\label{r15ray:distance}
Let $P\subset\mathbb R^N$ be a finite polyhedral set.  Let
$\lambda_{ij}=\lambda_{ji}:P\to\mathbb R$, $i\ne j$, be restrictions
of affine functions, satisfying
\[
 \lambda_{ij}\geq\min(\lambda_{ik},\lambda_{kj})
 \qquad(i,j,k\text{ distinct}).                 \tag{R15.1}
\]
Set $\lambda_{ii}=+\infty$.  For each $u\in P$ form the quotient
\[
 \mathcal T_u=(\{1,\ldots,m\}\times\mathbb R)/\sim,
 \qquad (i,t)\sim(j,t)\iff t\leq\lambda_{ij}(u).
                                                        \tag{R15.2}
\]
The parameter $t$ has weight one on each open edge.  The quotient embeds
as a polyhedral tree in $\mathbb R^m$ by
\[
 F_i(u,t)_j=\min(t,\lambda_{ij}(u));\qquad F_i(u,t)_i=t.
                                                        \tag{R15.3}
\]
Write $\mathcal T\subset P\times\mathbb R^m$ for this finite
polyhedral family and $F_i:P\times\mathbb R\to\mathcal T$ for the
maps including the $u$-coordinate.
\end{definition}

\begin{lemma}[The minimum-label decomposition]
\label{r15ray:labels}
For $i>1$ set $L_i(u)=\max_{j<i}\lambda_{ij}(u)$ and set
$L_1=-\infty$.  For a fixed $u$, the sets
\[
 F_1(u,\mathbb R),\qquad
 F_i(u,\{t>L_i(u)\})\quad(2\leq i\leq m)       \tag{R15.4}
\]
partition $\mathcal T_u$.  Using closed intervals instead changes only
finitely many vertices.  This partition retains the weight one on each
open edge.  The maps $F_i$ and the functions $L_i$ are finite piecewise
affine functions of the supplied distance coordinates, including at
equalities among the distances.
\end{lemma}

\begin{proof}
Condition (R15.1) proves transitivity in (R15.2).  If $t\leq
\lambda_{ij}$, the ultrametric condition implies
$\min(t,\lambda_{ik})=\min(t,\lambda_{jk})$ for every $k$;
thus the two expressions in (R15.3) agree.  Conversely equality of their
$i$ and $j$ coordinates implies $t\leq\lambda_{ij}$.  This proves the
embedding and the equivalence criterion.  Every equivalence class has
a unique least label $i$.  The representative $(i,t)$ has that label
exactly when $t>\lambda_{ij}$ for every $j<i$, which is exactly
$t>L_i(u)$.  This proves (R15.4).  There are only finitely many affine
comparisons; their common subdivision proves the piecewise affine
assertion.  No tree ordering has to be chosen afresh as $u$ moves.
\end{proof}

Here is the tensor version of interval integration needed below.  The
base may be nonpure-dimensional.  No weight or tropical-space structure
on that base is required.

\begin{lemma}[Integration on a moving band]
\label{r15ray:band}
Let $P\subset\mathbb R^N$ be a finite polyhedral set, let
$a,b:P\to\mathbb R$ be continuous piecewise affine functions, and let
$\theta\in\mathrm{PS}^{p,q}(P\times\mathbb R)$, $p,q\geq1$,
be given on a neighborhood of the bands in question.  Define
\[
 \mathcal I_{a,b}\theta
 =\int_{a(u)}^{b(u)}\operatorname{vc}_t(\theta)(u,t)\,dt
       \quad\text{if }a(u)<b(u),
 \qquad \mathcal I_{a,b}\theta=0\quad\text{if }a(u)\geq b(u).
                                                        \tag{R15.5}
\]
Here $\operatorname{vc}_t$ means the coefficient after contraction
with the two vertical tangent vectors, normalized by
$\operatorname{vc}_t(d't\wedge d''t\wedge\xi)=\xi$ for a base tensor
$\xi$.  Then (R15.5) is a PWS tensor of bidegree $(p-1,q-1)$ on $P$.
It commutes with restriction to every base face and with pullback by
a piecewise affine base map.  The assertion is local over relatively
compact subsets of $P$, so finite bounds on the bands suffice.
\end{lemma}

\begin{proof}
Refine $P$ so that $a,b$ are affine on its cells and add the cut $a=b$.
Intersect a finite complex of definition of $\theta$ with each band.
Projection of its finitely many faces gives a finite subdivision of the
base.  One-dimensional slicing over each resulting cell gives finitely
many intervals with affine endpoint functions.  On each slice the
contracted coefficients are smooth up to its closed faces.  For such a
coefficient $A$ the identity
\[
 \int_c^d A(u,t)\,dt
  =(d-c)\int_0^1 A(u,c+s(d-c))\,ds             \tag{R15.6}
\]
proves smoothness up to the base face, including when $c=d$.  The same
identity applies to every coefficient of the contracted tensor; the
finite number of smooth extensions supplies bounded coefficients on
compact parts.  This is also the affine fiber-integration construction
of Mihatsch, Section 2.3, equations (2.15)--(2.16).

For completeness, compatibility must be checked on a base face rather
than inferred from scalar coefficients.  If $\tau$ is a base face, both
vertical vectors are tangent to $\tau\times\mathbb R$.  Thus tangential
restriction commutes with their contraction.  Slices whose length tends
to zero have zero restricted tensor by (R15.6).  For every surviving
open slice, all limiting source presentations restrict to the same
tensor on $\tau\times\mathbb R$, by the PWS compatibility of $\theta$.
Its two vertical contractions therefore coincide.  The finitely many
slice endpoints have one-dimensional measure zero.  Summing proves
that all neighboring formulas restrict to (R15.5) on $\tau$.
The same argument applied to the graph of a piecewise affine base map
proves the stated base-change property.  This proves a PWS tensor on
the entire base, not merely smooth coefficients on its open cells.
No indicator function of the band is extended by zero as a tensor.
\end{proof}

\begin{proposition}[A finite formula on the distance family]
\label{r15ray:polyhedral}
Let $\Theta\in\mathrm{PS}^{p,q}(\mathcal T)$, $p,q\geq1$.
Suppose locally over a compact part of $P$ that the vertical contractions
of every $F_i^*\Theta$ vanish outside one interval $[T_-,T_+]$.
There is an intrinsic PWS tensor
$\mathcal R\Theta\in\mathrm{PS}^{p-1,q-1}(P)$ given by
\[
 \mathcal R\Theta
 =\sum_{i=1}^m
     \mathcal I_{\max(T_-,L_i),T_+}(F_i^*\Theta),          \tag{R15.7}
\]
where $\max(T_-,L_1)=T_-$.  It is independent of the bounds and of
the ordering of the labels.  It commutes with restrictions and
piecewise affine base refinements.  Only PWS pullback and band
integration occur in this construction.
\end{proposition}

\begin{proof}
Lemma~\ref{r15ray:band} gives a PWS tensor for every summand, proving
existence on the whole $P$.  Its base-change and restriction statements
give those for the sum.  Lemma~\ref{r15ray:labels} proves that the
formula is exactly integration once along each open fiber edge.
The contraction is invariant under a translated or reversed affine
fiber coordinate with its weight-one measure.  Consequently a second
label ordering gives the same edge integrals.  Vertices have
one-dimensional measure zero; collapsed intervals contribute zero by
(R15.6).  Bounds outside the vertical support do not change the
integrals.  These observations also prove compatibility when a
combinatorial edge disappears or several distance functions coincide.
\end{proof}

The last assertion does not claim that scalar coefficients of an
arbitrary PWS tensor are continuous across its faces.  It is the
contracted tensor, with its tangential restriction, that is integrated.
For instance, a vertical coefficient may jump across $t=u$ without
violating either lemma.

\begin{proposition}[Ordinary realization and all smooth mixed tests]
\label{r15ray:analytic}
Let $B$ be an ordinary boundaryless analytic open of pure dimension $d$
over the fixed field $K$.  Let $a_1=0,a_2,\ldots,a_m\in\mathcal O(B)$
be pairwise distinct analytic sections such that every difference
$a_i-a_j$ is a unit.  Choose a finite tuple of analytic units $b$ on
$B$ containing these nonzero differences and all other base units
used below.  Put
\[
 W=(B\times\mathbb A^1_z)\setminus
                 \bigcup_i\{z=a_i\},\qquad w_i=z-a_i.
\]
Let $\theta\in\mathrm{PS}^{p,q}(W)$ have one finite PWS presentation
in the genuine unit tuple $(b,w_1,\ldots,w_m)$.  Assume its support
is proper over $B$ and, locally over each relatively compact target
chart, is contained away from uniform tails at all marked sections
and at infinity.  Then the current pushforward along $p:W\to B$
is represented by an ordinary actual PWS tensor of degree
$(p-1,q-1)$, for $p,q\geq1$.  More precisely the local representative
is the pullback of (R15.7) through the finite unit tuple $b$, and
\[
 \int_B\mathcal R_B(\theta)\wedge\eta
     =\int_W\theta\wedge p^*\eta                      \tag{R15.8}
\]
for every ordinary smooth compactly supported complementary test
$\eta$, without an exterior-product restriction.  If $p=0$ or $q=0$,
the corresponding pushforward current is zero by relative degree.
\end{proposition}

\begin{proof}
\emph{A genuine ordinary candidate.}
Work on a relatively compact affinoid neighborhood in $B$ whose
interior is an ordinary neighborhood of the point under consideration.
Its $b$-image is a finite polyhedral set $P$.  The distance functions
$\lambda_{ij}=v(a_i-a_j)$ are coordinate functions on $P$ and satisfy
(R15.1).  For each actual base point with image $u$, the Gauss
valuation of the disk centered at $a_i$ with valuation radius $t$
has $v(w_j)=\min(t,\lambda_{ij})$.  Thus the whole family
$\mathcal T$ maps into the actual moment image of $W$.
Restrict the given PWS presentation to this family, obtaining
$\Theta$.  The tail hypothesis and boundedness of the distances on
the compact target chart give finite uniform $T_-,T_+$.  Formula
(R15.7) gives a PWS tensor on $P$.  The finite-polyhedron PWS
extension lemma extends it to a neighborhood of $P$ in its ambient
Euclidean space.  Pullback through the actual units $b$ therefore
gives an ordinary PWS, hence actual, tensor.  This last step is why
an intrinsic PWS output on an unspecified skeleton would not suffice.

\emph{The skeleton after adding all test moments.}
Take an ordinary chart for the arbitrary test $\eta$ and append all
its unit moments to $b$, obtaining $b'$.  Write
$S=\Sigma(B,b')$ on this chart.  The refined source skeleton for
$(b',w_1,\ldots,w_m)$ is the marked-tree family over $S$.
Here is a direct verification, including the projection to $S$.
At any source point choose $i$ for which $v(w_i)$ is maximal.
Then for each $j$ the ultrametric rule gives
$v(w_j)=\min(v(w_i),v(a_i-a_j))$.  At equality a cancellation
would give $v(w_j)>v(w_i)$, contradicting this choice.  Hence every
graded reduction of $w_j$ is generated by that of $w_i$ and by
the graded reductions of the differences, all included in $b'$.
The source graded transcendence degree is therefore at most the
base graded transcendence degree plus one.  A source point of
the $(d+1)$-dimensional skeleton projects to $S$.  Over such a
base point the graded field generated by $b'$ has transcendence
degree $d$, so the full graded residue field of the base is algebraic
over it.  The additional initial form of $w_i$ is therefore
transcendental over the full graded residue field as well.  In a
polynomial in $w_i$ this prevents cancellation among all terms of
minimal valuation.  It gives the Gauss norm formula, and thus says
precisely that the point is a Gauss point in the coordinate $w_i$.
Conversely these Gauss points have the required transcendence
degree.  Their valuations are exactly (R15.3).

On an open edge, $w_i=z-a_i$ is a genuine analytic coordinate
over the base.  The remaining moment logarithms are affine in
the base logarithms and $t=v(w_i)$ there.  The coordinate-product
skeleton lemma of Round 11 identifies this edge with the
corresponding product band in $S\times\mathbb R$, with the native
base weight times the standard weight one.  This statement is
compatible with every appended test tuple.  The argument does
not use the dimension of the initial coarse image $P$.

\emph{Integration.}
The paired input $\theta\wedge p^*\eta$ is PWS of top degree
$(d+1,d+1)$.  Its support lies on the refined skeleton, by
Mihatsch Corollary 4.9(2).  Integration is the intrinsic skeletal
integration of Definition 4.10.  The relative vertices lie on finitely
many graphs $t=\lambda_{ij}$ over $S$ and have dimension at most
$d$.  They carry no PWS top-degree measure.  On every remaining
open edge the product weight just computed gives ordinary
coefficientwise Euclidean Fubini after contraction in both
vertical directions.  Restricting the resulting top measure
to a band is legitimate; no PWS zero extension across its
boundary is made.  The finite number of coefficient pieces,
their boundedness on compact paired support, and the uniform
tail bounds prove absolute integrability.  Summing the
minimum-label bands gives (R15.7), pulled back to every test
refinement.  Its base-change property gives the paired top
form $\mathcal R_B(\theta)\wedge\eta$ on $S$.
Intrinsic skeletal integration proves (R15.8) on this chart.
Ordinary smooth partitions of unity on the target prove it
for every compactly supported $\eta$.

Any two constructed ordinary representatives thus have the same
current on their ordinary overlap.  Mihatsch Proposition 4.11
makes them equal as actual forms.  They glue in the ordinary
topology and are independent of all presentations, bounds,
extensions, and label orderings.  Native weights have been
retained throughout.  No auxiliary field extension occurs in
this proposition itself.
\end{proof}

\paragraph{Explicit low-degree integrands.}
On one band write the input as
\[
 \alpha=A\,d't\wedge d''t
   +\sum_j B_j\,d'u_j\wedge d''t
   +\sum_j C_j\,d't\wedge d''u_j+\alpha_{\mathrm{hor}}.
\]
For a scalar kernel $g$ and a common scalar cutoff $\chi$, the
band contributions are
\[
 G_{\mathrm{band}}=\int\chi gA\,dt,\qquad
 P_{\mathrm{band}}
   =\int\chi\left(A\,d'_u g-g_t\sum_j B_jd'u_j\right)\,dt.
                                                        \tag{R15.9}
\]
The second term in $P$ is essential: scalar edge-length integration
does not recover it.  Derivatives of $\chi$ are absent because
$P$ integrates $\chi\,d'g\wedge\alpha$, not
$d'(\chi g)\wedge\alpha$.  The differential identity for the
operators must be derived before cutoff insertion, as in Round 14.

\paragraph{Dependencies and limitations.}
The public inputs are PWS calculus, actual current injectivity,
support on the presenting skeleton, and native integration in
Mihatsch, \emph{$\delta$-Forms on Lubin--Tate spaces}, Sections
2.2--2.4, 3.2--3.3, and 4.1--4.3.  The product weight and arbitrary
test-refinement assertion are the separately audited Round 11
coordinate-product skeleton lemma.  The finite-polyhedron extension
is the Round 14 PWS extension lemma; it is degreewise and is not a
closed ambient lifting theorem.  The verification of the global
Green integrands against the hypotheses above, including collision
charts and the bad base fibers, must be supplied separately.
